% ==============================================================================
% MynaNet: Efficient Audio-First CNN for Edge-Based Avian Acoustic Classification
% Target journal: Neurocomputing (Elsevier)
% Format: elsarticle
% ==============================================================================

%\documentclass[review,12pt]{elsarticle}
\documentclass[final,5p,times,twocolumn,nolinenum]{elsarticle}

\usepackage{amsmath,amssymb}
\usepackage{graphicx}
\usepackage{booktabs}
\usepackage{multirow}
\usepackage{hyperref}
\usepackage{lineno}
\usepackage{xcolor}
\usepackage{microtype}
\usepackage{siunitx}

% Convenience command for TODO annotations
\newcommand{\TODO}[1]{\textcolor{red}{\textbf{[TODO: #1]}}}

\journal{Neurocomputing}

\begin{document}

\begin{frontmatter}

\title{MynaNet: A Lightweight Audio-First Convolutional Neural Network for
       Edge-Based Avian Acoustic Classification}

%% Authors
%\author[um]{Muhammad Mun'im Ahmad Zabidi\corref{cor1}}
%\ead{mzabidi@gmail.com}
%\author[um]{Muhammad Yamani Idna Idris}
%\author[um]{Norisma Idris}
\cortext[cor1]{Corresponding author}

%\address[um]{Universiti Malaya, Kuala Lumpur, Malaysia}

% ==============================================================================
\begin{abstract}
We present MynaNet, a lightweight audio-first convolutional neural
network designed for automated bird call classification on
resource-constrained edge devices. Through systematic architectural
exploration spanning computer-vision transfer-learning baselines
(MobileNetV3), temporal sequence models (Temporal Convolutional Networks),
and audio-specific depthwise-separable CNNs with squeeze-and-excitation
attention, we demonstrate that domain-specific inductive biases
substantially outperform generic architectures within severe memory
budgets. MynaNet~v1 achieves \num[round-mode=places,round-precision=2]{95.00}\%
$\pm$ \SI{0.87}{\percent} mean INT8 accuracy at a \SI{434}{KB} model size
on ten Southeast Asian garden-bird species, validating the deployment
targets of $>$\SI{90}{\percent} accuracy within $<$\SI{512}{KB}. Ablation
studies across 20+ configurations isolate four primary design insights:
(i)~audio-aware 2D convolutions outperform 1D temporal models by
\SI{2.2}{\percent} absolute; (ii)~Mixup augmentation surpasses SpecAugment
by \SI{0.5}{\percent}; (iii)~moderate multi-head self-attention (2~heads,
88 dimensions) provides the optimal accuracy-size trade-off; and
(iv)~a 6\%-reduced channel width ([120,~180,~146,~360]) improves
generalisation over a wider over-parameterised baseline. The resulting
architecture is \SI{15}{\percent} under the \SI{512}{KB} deployment budget
with an estimated latency of \SI{116}{ms} on an ARM Cortex-M7 at
\SI{480}{MHz}, enabling continuous real-time acoustic monitoring for
wildlife conservation.

% TODO: Update abstract with final 3-seed results on MyGardenBird finalised
%       dataset (80:10:10 split) after Linux retraining.  Current authoritative
%       figures (95.00 +- 0.87%) are from OLD dataset (Common Myna + Zebra
%       Dove, 80:10:10 split).  macOS seed-42 preliminary on NEW dataset
%       (Coppersmith Barbet + White-breasted Waterhen, 80:10:10): pending;
%       prior 75:10:15 run gave 93.44% FP32 no-mixup / 93.89% FP32 mixup /
%       94.00% INT8.
\end{abstract}

\begin{keyword}
bird call classification \sep
edge AI \sep
depthwise-separable CNN \sep
squeeze-and-excitation \sep
multi-head self-attention \sep
post-training quantisation \sep
wildlife acoustics \sep
embedded systems
\end{keyword}

\end{frontmatter}

\linenumbers

% ==============================================================================
\section{Introduction}
\label{sec:introduction}
% ==============================================================================

Automated acoustic monitoring of bird populations is a cornerstone of
biodiversity surveying and ecological conservation research. Deploying
classifiers directly on low-power acoustic sensors eliminates the cost and
latency of cloud uplinks, enabling dense, long-term monitoring at scale.
However, meeting accuracy requirements within the extreme memory constraints
of ARM Cortex-class microcontrollers ($<$\SI{512}{KB} flash) demands
purpose-built model architectures rather than off-the-shelf computer-vision
backbones.

This paper presents MynaNet~v1, the outcome of a principled multi-phase
architecture search starting from pretrained MobileNetV3 baselines, passing
through a temporal-convolutional-network detour, and arriving at a
hand-optimised audio-first design. The dataset used for all experiments is
the MyGardenBird corpus~\cite{zabidi2026mygardenbird}, a curated twelve-class
collection of Southeast Asian garden-bird recordings sourced from Xeno-canto
and segmented into three-second clips.

The main contributions of this work are:
\begin{enumerate}
  \item A systematic ablation study (Section~\ref{sec:ablation}) documenting
        the incremental gains from depthwise-separable convolutions,
        squeeze-and-excitation channel attention, residual skip connections
        in all blocks, and lightweight multi-head self-attention.
  \item MynaNet~v1 (Section~\ref{sec:architecture}), which achieves 95.00\%
        INT8 accuracy at 434~KB on the 10-class dataset and 93.52\% on the
        expanded 12-class dataset (Section~\ref{sec:reproducibility})---\SI{38}{\percent}
        smaller than a comparable MobileNetV3-Small variant at equivalent or superior
        accuracy.
  \item Quantitative evidence that 2D audio-specific inductive bias
        outperforms 1D temporal models by \SI{2.2}{\percent} despite equal
        parameter budgets (Section~\ref{sec:tcn}).
  \item An INT8-quantised TFLite deployment profile meeting all edge
        constraints (Section~\ref{sec:deployment}).
\end{enumerate}


% ==============================================================================
\label{sec:related}
% ==============================================================================
\section{Related Work}

\subsection{Bioacoustic Classification Systems}


\paragraph{Narrative point}:
Large server-class models dominate the field and achieve high accuracy,
but none target MCU deployment. Establish the importance of the task and
the server-class "ceiling" that MynaNet must approach from below.

Automated bird-call classification has advanced significantly with the availability of large-scale citizen-science datasets such as Xeno-canto and eBird. Modern systems typically employ deep convolutional neural networks (CNNs) trained on log-mel spectrogram representations. One of the most widely deployed systems is BirdNET \cite{kahl2021birdnet}, which combines convolutional architectures with curated global training data to achieve large-scale multi-species recognition. Integrated into the Merlin Sound ID application, BirdNET demonstrates that deep learning can approach expert-level identification performance in real-world acoustic monitoring scenarios.

Large-scale pretrained audio models have also been applied to bioacoustics. PANNs (Pretrained Audio Neural Networks) \cite{kong2020panns} leverage CNN architectures trained on AudioSet to enable transfer learning for downstream acoustic classification tasks. In competitive benchmarking environments such as the BirdCLEF challenges \cite{goerlich2022overview}, top-performing systems frequently adopt EfficientNet or ResNet-style backbones trained with extensive augmentation and ensembling strategies.

While these approaches establish state-of-the-art accuracy, they typically rely on GPU-class hardware and models containing millions of parameters. Model size, energy consumption, and flash memory constraints are rarely primary considerations in these architectures. Consequently, despite their effectiveness, such systems are not directly suitable for continuous deployment on microcontroller-class edge devices.

\subsection{Lightweight Convolutional Networks for Audio}

Efficient convolutional design has enabled neural networks to operate under constrained computational budgets. Depthwise-separable convolutions, popularized by MobileNet \cite{howard2017mobilenets}, substantially reduce parameter counts and multiply-accumulate (MAC) operations relative to standard convolutions. MobileNetV3 further refines this paradigm using neural architecture search (NAS) and lightweight attention mechanisms \cite{howard2019searching}.

In the audio domain, depthwise-separable CNN (DS-CNN) architectures for keyword spotting (KWS) demonstrated that compact models can achieve competitive accuracy while remaining deployable on microcontrollers \cite{zhang2017hello}. These models exploit the 2D structure of log-mel spectrograms while maintaining strict computational efficiency. EfficientNet variants \cite{tan2019efficientnet} have also been adapted to spectrogram-based classification tasks, though such models typically exceed the SRAM and flash budgets of low-power embedded systems.

Most lightweight CNN research in audio focuses on speech recognition or KWS. Wildlife acoustics presents distinct challenges, including harmonic stacking, frequency sweeps, and species-specific trill structures, which may benefit from task-specific architectural adaptation rather than direct porting of speech models.

\subsection{Edge AI and TinyML for Acoustic Monitoring}

The TinyML paradigm seeks to deploy neural networks on microcontrollers with flash budgets typically below 1~MB. Frameworks such as TensorFlow Lite Micro \cite{david2021tensorflow} and Arm CMSIS-NN \cite{lai2018cmsis} provide optimized INT8 inference kernels for Cortex-M processors. MCUNet \cite{lin2020mcunet} further demonstrates that the co-design of neural architectures and inference runtimes enables sub-megabyte deep learning models for embedded systems.

Within acoustic applications, TinyML deployments have predominantly targeted KWS. Although environmental sound monitoring has been explored, few published studies report bird-call classifiers that simultaneously achieve high multi-class accuracy and verified sub-512KB flash deployment using quantized INT8 inference. Thus, a gap remains between high-accuracy bioacoustic models and strictly memory-constrained embedded deployments.

\subsection{Attention Mechanisms in Audio Models}

Attention mechanisms have become central to modern audio modeling. The Audio Spectrogram Transformer (AST) \cite{gong2021ast} demonstrates that full self-attention over spectrogram patches can outperform convolutional baselines in large-data regimes. Hybrid CNN--attention architectures have also been proposed for acoustic event detection and speech modeling \cite{kong2020panns}.

Channel-wise recalibration via Squeeze-and-Excitation (SE) blocks improves representational efficiency with minimal overhead \cite{hu2018squeeze}. However, transformer-style architectures generally incur substantial computational and memory costs, limiting their direct applicability to microcontroller-class hardware. Lightweight multi-head self-attention integrated within compact CNN backbones remains comparatively underexplored in embedded bioacoustic classification scenarios.

\subsection{Data Augmentation for Acoustic Classification}

Data augmentation is critical for improving generalization in audio classification. Mixup \cite{zhang2018mixup} linearly interpolates training examples and labels, improving calibration and robustness. SpecAugment \cite{park2019specaugment} applies time and frequency masking directly to spectrogram representations and has demonstrated strong performance in speech recognition tasks. While both techniques are widely adopted, their relative effectiveness in small-scale wildlife acoustic datasets remains insufficiently characterized, particularly under constrained training data regimes.

\subsection{Positioning of This Work}

Existing literature demonstrates three largely parallel trajectories: (i) high-accuracy bird-call classifiers built on large CNN or transformer backbones; (ii) lightweight CNN architectures primarily evaluated on speech tasks; and (iii) TinyML deployment frameworks targeting microcontroller inference. However, few studies unify these directions into a controlled architectural exploration for bird-call classification under strict sub-512KB flash constraints.

To our knowledge, no prior work systematically compares pretrained vision backbones, 1D temporal convolutional networks, and audio-specific 2D depthwise-separable CNNs within a fixed microcontroller deployment budget for avian acoustic recognition. MynaNet addresses this gap through a multi-phase architectural investigation culminating in a 434KB INT8 model validated via controlled ablations and multi-seed evaluation.

% ==============================================================================
\section{Dataset: MyGardenBird}
\label{sec:dataset}
% ==============================================================================

All experiments use the MyGardenBird corpus~\cite{zabidi2026seabird}, a
twelve-class dataset of Southeast Asian garden birds recorded in Singapore and
peninsular Malaysia. Audio is sourced from Xeno-canto and segmented into
non-overlapping three-second clips at \SI{16}{kHz}. The twelve species are:
Asian Koel, Collared Kingfisher, Common Iora, Common Tailorbird, Coppersmith
Barbet, Large-tailed Nightjar, Olive-backed Sunbird, Pied Fantail, Spotted Dove,
White-breasted Waterhen, White-throated Kingfisher, and Yellow-vented Bulbul.
Each species contributes 600 clips (7{,}200 total).

All experiments use the 80:10:10 train/val/test split, providing
5{,}760 training, 720 validation, and 720 test samples (60 per species
per partition). This split maximises training data at the 7{,}200-clip
scale while retaining a held-out test set sufficient for unbiased
evaluation. The split was generated using Mixed Integer Programming to
ensure source-based separation and prevent data leakage across recordings.
The split files are distributed with the corpus under \textsc{cc-by}~4.0.

\TODO{Add dataset provenance paragraph: number of Xeno-canto source
recordings per species, total clip counts before/after quality filtering,
and cite zabidi2026seabird once the Scientific Data paper is accepted.}

% ==============================================================================
\section{Feature Extraction}
\label{sec:features}
% ==============================================================================

All models consume log-mel spectrograms computed with
Librosa~\cite{TODO-librosa}:
\begin{itemize}
  \item Sampling rate: \SI{48}{kHz}; clip length: 3 seconds
  \item FFT window: 2{,}048 samples (\SI{42.7}{ms})
  \item Hop length: 160 samples (\SI{10}{ms stride}) $\Rightarrow$
        $T = 48{,}000 / 160 = 300$ frames per clip
  \item Mel bins: $M = 80$ (0--\SI{8}{kHz}, Slaney-scale filter bank)
  \item Output: $80 \times 300$ log-power spectrum (dB, no $\text{top\_db}$
        clamp), normalised to $[0, 255]$ and replicated across 3 channels
        for compatibility with pretrained visual backbones
\end{itemize}

The 80-bin configuration provides \SI{100}{Hz/bin} spectral resolution,
sufficient to resolve harmonics of the high-frequency species in the dataset
(e.g., Olive-backed Sunbird, 4--\SI{7}{kHz}).

% ==============================================================================
\section{Architecture Design and Evolution}
\label{sec:architecture}
% ==============================================================================

We chronicle the architectural progression across four phases below.

% ------------------------------------------------------------------------------
\subsection{Phase 1: Computer Vision Baselines (MobileNetV3)}
\label{sec:mobilenet}
% ------------------------------------------------------------------------------

We began with MobileNetV3-Small~\cite{howard2019searching} pretrained on
ImageNet-1k, adapting it for $80 \times 300$ mel-spectrogram
classification. Width multipliers were swept to find the accuracy-size
frontier (Table~\ref{tab:mobilenet_baseline}).

\begin{table}[ht]
\centering
\caption{MobileNetV3-Small results.
         Upper block: width-multiplier sweep, $80\times300$ input, 80:10:10 split,
         OLD dataset (Common Myna + Zebra Dove).
         Lower block: native $80\times300$ audio-input baseline,
         NEW finalised dataset (Coppersmith Barbet + White-breasted Waterhen),
         FP32 test accuracy, 3-seed mean $\pm$ s.d.\ (macOS; $^\ddagger$denotes
         macOS platform, not Linux authoritative).}
\label{tab:mobilenet_baseline}
\smallskip\small
\resizebox{\linewidth}{!}{%
\begin{tabular}{@{}lrrrr@{}}
\toprule
\textbf{Variant} & \textbf{Params} & \textbf{Size (KB)} &
\textbf{FP32 Acc.} & \textbf{INT8 Acc.} \\
\midrule
\multicolumn{5}{l}{\textit{Width sweep --- OLD dataset, 80:10:10}} \\
Width $\times$1.0 & 2.54M & 2{,}478 & 94.44\% & 93.56\% \\
Width $\times$0.75 & 1.53M & 1{,}495 & 94.11\% & 93.11\% \\
Width $\times$0.5  & 0.72M & 703     & 92.78\% & 91.67\% \\
Width $\times$0.35 & 0.39M & 381     & 90.44\% & 89.11\% \\
\midrule
\multicolumn{5}{l}{\textit{Native 80$\times$300 input --- NEW dataset, 80:10:10}} \\
Width $\times$1.0 (Stage9b)$^\ddagger$ & 1.24M & \textasciitilde1{,}200 &
  $82.22 \pm 2.24$\% & -- \\
\bottomrule
\multicolumn{5}{l}{\footnotesize $^\ddagger$3-seed macOS (seed 42: 83.17\%,
  seed 100: 79.67\%, seed 786: 83.83\%).} \\
\end{tabular}}
\end{table}

All width variants exceed the \SI{512}{KB} deployment constraint, and none
is optimised for spectro-temporal patterns. On the finalised dataset with
native $80\times300$ input, MobileNetV3-Small achieves only
$82.22 \pm 2.24$\% (FP32, 3-seed macOS), compared to MynaNet~v1's
93.44\% FP32 on the same platform and split (Section~\ref{sec:results}),
a gap of \SI{11.22}{\percent} absolute. MobileNetV3's isotropic 2D
convolutions treat frequency and time equivalently, and its early stride-2
operations discard fine-grained temporal structure. These limitations
motivated an audio-specific design.

% ------------------------------------------------------------------------------
\subsection{Phase 2: Audio-Specific Depthwise-Separable CNNs}
\label{sec:dscnn}
% ------------------------------------------------------------------------------

Inspired by the DS-CNN architecture for keyword spotting~\cite{zhang2017hello},
we designed a four-block depthwise-separable backbone tailored to
spectrograms. Table~\ref{tab:dscnn_evolution} documents the incremental
ablation from the baseline (Model~1a) to the final DS configuration
(Model~1e).

\begin{table}[ht]
\centering
\caption{DS-CNN architecture evolution.
         Input: $80\times300$, 80:10:10 split, INT8 accuracy.
         Results on OLD dataset.}
\label{tab:dscnn_evolution}
\resizebox{\linewidth}{!}{%
\begin{tabular}{@{}llrrl@{}}
\toprule
\textbf{Model} & \textbf{Key Innovation} &
\textbf{Params} & \textbf{Size (KB)} & \textbf{INT8 Acc.} \\
\midrule
1a: DS-CNN Baseline  & Depthwise-separable blocks   & 194K & 189 & 91.67\% \\
1b: +Wider Stem      & 64-channel initial conv       & 297K & 290 & 92.89\% \\
1c: +SE Attention    & Squeeze-and-Excitation        & 318K & 311 & 93.56\% \\
1d: +Residual+MHSA   & Skip connections + attention  & 309K & 302 & 94.22\% \\
\textbf{1e: +All Residuals} & All 4 blocks + 80 mels & \textbf{462K} & \textbf{451} & \textbf{95.67\%} \\
\midrule
v0: Production Baseline & Channels [128,192,156,384] & 395K & 481 & 93.00\% \\
\textit{v1: MynaNet (Final)} & \textit{6\% reduction [120,180,146,360]} &
  \textit{328K} & \textit{434} & \textit{\textbf{95.00\%$^\dagger$}} \\
v1sa: +SpecAugment & Frequency/time masking         & 395K & 434 & 94.17\% \\
v2: Enhanced MHSA  & 3 heads, 112 dims              & 372K & 477 & 94.33\% \\
\bottomrule
\multicolumn{5}{l}{\footnotesize $^\dagger$Mean of 3-seed Linux evaluation
  (94.00\%, 95.50\%, 95.50\%; $\sigma=0.87$\%).} \\
\end{tabular}}
\end{table}

Progressive improvements are as follows.
\begin{itemize}
  \item \textbf{1a $\to$ 1b (+1.22\%):} A wider 64-channel stem captures
        richer initial spectro-temporal features.
  \item \textbf{1b $\to$ 1c (+0.67\%):} Squeeze-and-Excitation
        attention~\cite{hu2018squeeze} recalibrates channel responses,
        suppressing noise-dominated bands.
  \item \textbf{1c $\to$ 1d (+0.66\%):} Skip connections improve gradient
        flow; MHSA captures long-range temporal dependencies such as trill
        patterns.
  \item \textbf{1d $\to$ 1e (+1.45\%):} Extending residuals to all four
        blocks and increasing to 80~mel bins enables \SI{90}{epoch} training
        without degradation.
  \item \textbf{1e $\to$ v1 (+0.33\% mean, $-$95~KB):} A strategic 6\%
        channel reduction eliminates over-parameterisation; 3-seed Linux
        evaluation confirms v1 (95.00\%) surpasses the larger 1e
        (94.67\%~$\pm$~0.45\%).
\end{itemize}

% ------------------------------------------------------------------------------
\subsection{Phase 3: Temporal Convolutional Networks}
\label{sec:tcn}
% ------------------------------------------------------------------------------

Motivated by the empirical superiority of TCNs over recurrent networks for
sequence modelling~\cite{bai2018empirical}, we explored whether causal
dilated 1D convolutions could match our 2D DS-CNN while operating on
pre-pooled spectrogram embeddings. Table~\ref{tab:tcn_ablation} reports
eight controlled TCN variants.

\begin{table}[ht]
\centering
\caption{TCN ablation study.
         Input: $64\times300$ embeddings, Mixup $\alpha=0.2$, 80:10:10 split.
         Results on OLD dataset.}
\label{tab:tcn_ablation}
\resizebox{\linewidth}{!}{%
\begin{tabular}{@{}llrrl@{}}
\toprule
\textbf{Variant} & \textbf{Modification} &
\textbf{Params} & \textbf{Size (KB)} & \textbf{INT8 Acc.} \\
\midrule
4a: Baseline (8 layers) & Dilation [1,2,4,8,16,32,64,128] & 242K & 236 & 92.00\% \\
4b: Shallow (4 layers)  & Dilation [1,2,4,8]              & 162K & 158 & 91.33\% \\
4c: Wide (128 ch)       & $2\times$ channel width          & 461K & 450 & 92.44\% \\
4d: Deep (10 layers)    & +2 layers, dilation to 512       & 298K & 291 & 92.22\% \\
4e: Kernel-2            & $k=2$ (RF = 255 frames)          & 194K & 189 & 91.56\% \\
4f: Kernel-5            & $k=5$ (RF = 1{,}023 frames)     & 337K & 329 & 92.11\% \\
4g: No-residual         & Removed skip connections          & 242K & 236 & 90.89\% \\
4h: Lightweight         & 32 channels, 6 layers            & 142K & 139 & 90.67\% \\
\midrule
\textit{DS-CNN 1e (ours)} & \textit{2D hybrid+attention} &
  \textit{462K} & \textit{451} & \textit{\textbf{95.67\%}} \\
\bottomrule
\end{tabular}}
\end{table}

The best TCN variant (4c, 92.44\%) lags DS-CNN~1e by 3.23\% absolute.
Three factors explain this gap. First, TCN operates on 1D embeddings
produced by aggressive 2D pooling, collapsing spectral structure that
discriminates species by harmonic stacks. Second, the 511-frame receptive
field (5.1~s) far exceeds the 3-second clip length, wasting capacity on
zero-padding artifacts. Third, classification is non-causal (the entire
clip is available at inference), so causal padding incurs computation
without benefit. We conclude that spectrogram classification benefits from
2D spatial inductive bias over 1D temporal modelling.

% ------------------------------------------------------------------------------
\subsection{Phase 4: MynaNet v1 --- Final Architecture}
\label{sec:mynanet}
% ------------------------------------------------------------------------------

Figure~\ref{fig:architecture} illustrates MynaNet~v1. The network comprises:
\begin{enumerate}
  \item \textbf{Stem:} Conv2D(128, $3\times3$) + BatchNorm + ReLU6.
  \item \textbf{Four DS-Conv-SE blocks} with residual connections and
        $2\times2$ max-pooling, channel widths [120, 180, 146, 360].
        Blocks 2--4 use $1\times1$ projection in the skip path when channel
        counts change.
  \item \textbf{Lightweight MHSA:} Projects 360-dim features to 88-dim;
        2 heads with $d_k=32$; attention over $5\times18=90$ patches.
        LayerNorm applied after residual add.
  \item \textbf{Head:} GlobalAveragePooling $\to$ Dense(176) + BatchNorm +
        ReLU6 + Dropout(0.05) $\to$ Dense(10) + Softmax.
\end{enumerate}

\begin{figure*}[ht]
\centering
\includegraphics[width=\linewidth]{architecture.png}
%\fbox{\parbox{0.9\linewidth}{\centering\vspace{2cm}
%  [MynaNet v1 architecture diagram --- INSERT FIGURE]
%  \vspace{2cm}}}
\caption{MynaNet v1 architecture: $80\times300$ log-mel spectrogram input,
         four depthwise-separable blocks with SE attention and residual
         connections, lightweight 2-head self-attention, and dense
         classification head. Total: 327{,}908 parameters, 434~KB INT8.}
\label{fig:architecture}
\end{figure*}

\paragraph{SE reduction ratios.}
The squeeze-and-excitation bottleneck ratio scales adaptively: Block~1
(120~ch, $r=8$), Block~2 (180~ch, $r=12$), Blocks~3--4 (146--360~ch,
$r=12$--16), maintaining SE overhead at 1--2\% of block parameters.

\paragraph{Dropout schedule.}
Gradual dropout [0.025, 0.025, 0.0375, 0.05] across blocks matches
receptive field growth: early blocks (local context) tolerate less dropout;
final blocks (global context) require stronger regularisation.

% ==============================================================================
\section{Experimental Setup}
\label{sec:setup}
% ==============================================================================

\subsection{Training Configuration}

All models use two-phase training:
\begin{enumerate}
  \item \textbf{Warmup} (70~epochs): Adam / Legacy-Adam at $10^{-3}$ with
        cosine decay. Full network trainable.
  \item \textbf{Fine-tune} (20~epochs): Learning rate $10^{-5}$, refines
        decision boundaries.
\end{enumerate}

Mixup~\cite{zhang2018mixup} augmentation with $\alpha=0.2$ is applied
during warmup. On Apple Silicon, TensorFlow legacy Adam provides
$\approx$20\% speedup over AdamW with identical accuracy.

Post-training quantisation (PTQ) to INT8 is performed with 200 calibration
samples from the validation set. Quantisation incurs $<$\SI{0.5}{\percent}
accuracy drop for MynaNet~v1.

\subsection{Evaluation Protocol}

Ablation experiments (Section~\ref{sec:ablation}) use an 80:10:10 split
to maximise training data (5{,}760 / 720 / 720 samples). Final
reproducibility results (Section~\ref{sec:reproducibility}) use three
random seeds (42, 100, 786) on macOS (Apple M4 Pro), with Linux/CUDA
re-runs planned for authoritative platform-independent figures.

% ==============================================================================
\section{Results}
\label{sec:results}
% ==============================================================================

\subsection{Architecture Ablation}
\label{sec:ablation}

Table~\ref{tab:dscnn_evolution} (Section~\ref{sec:dscnn}) documents the
incremental gains of each architectural component. Extending skip
connections to all four DS blocks (1d $\to$ 1e) delivers the largest single
gain (+1.45\%), followed by the wider stem (1a $\to$ 1b, +1.22\%).
Lightweight MHSA adds only 13K parameters but contributes +0.66\% (1c
$\to$ 1d), confirming that long-range temporal attention complements local
depthwise convolutions. Scaling to 3 heads (v2) provides no benefit over
2 heads at higher cost, indicating attention-capacity saturation.

\subsection{Training Hyperparameter Ablation}

Table~\ref{tab:training_ablation} isolates the effect of augmentation
strategy, data split ratio, and attention capacity.

\begin{table}[ht]
\centering
\caption{Training hyperparameter ablation (MynaNet~v1, 80~mels, INT8 accuracy).
         Results on OLD dataset.}
\label{tab:training_ablation}
\resizebox{\linewidth}{!}{%
\begin{tabular}{@{}llccc@{}}
\toprule
\textbf{Group} & \textbf{Config} & \textbf{Augmentation} &
\textbf{Split} & \textbf{INT8 Acc.} \\
\midrule
\multirow{2}{*}{Augmentation} &
  Mixup $\alpha=0.2$    & Mixup 0.2       & 80:10:10 & \textbf{95.00\%$^\dagger$} \\
& SpecAugment            & Freq/Time Mask  & 80:10:10 & 94.17\% \\
\midrule
\multirow{3}{*}{Split ratio} &
  80:10:10               & Mixup 0.2       & 80:10:10 & \textbf{95.00\%$^\dagger$} \\
& 75:10:15               & Mixup 0.2       & 75:10:15 & 94.00\% \\
& 70:15:15               & Mixup 0.2       & 70:15:15 & 93.44\% \\
\midrule
\multirow{3}{*}{Architecture} &
  v1 (Optimized)         & Mixup 0.2       & 80:10:10 & \textbf{95.00\%$^\dagger$} \\
& v2 (MHSA 3h/112d)      & Mixup 0.2       & 80:10:10 & 94.33\% \\
& v0 (Baseline)          & Mixup 0.2       & 80:10:10 & 93.00\% \\
\bottomrule
\multicolumn{5}{l}{\footnotesize $^\dagger$3-seed Linux mean (seed 42:
  94.00\%, seed 100: 95.50\%, seed 786: 95.50\%).} \\
\end{tabular}}
\end{table}

Mixup outperforms SpecAugment by \SI{0.83}{\percent} on this dataset.
With only 4{,}800 training samples (10-class dataset), SpecAugment's aggressive masking risks
erasing discriminative short vocalisations ($<$\SI{200}{ms}), whereas
Mixup's label smoothing improves calibration without destructive masking.
Reducing training data from 80:10:10 (5{,}760) to 70:15:15 (5{,}040)
degrades accuracy by 1.56\%, underscoring the importance of maximising
training set size at this dataset scale.

\subsection{Comparison with Vision Baselines}
\label{sec:vision_comparison}

Table~\ref{tab:final_comparison} summarises the full architecture evolution.

\begin{table}[ht]
\centering
\caption{Architecture comparison across phases.
         $^\dagger$3-seed Linux mean (10-class dataset, 80:10:10).
         $^\ddagger$3-seed macOS mean $\pm$ s.d.\ (12-class dataset, 80:10:10, FP32).}
\label{tab:final_comparison}
\resizebox{\linewidth}{!}{%
\begin{tabular}{@{}llrrl@{}}
\toprule
\textbf{Phase} & \textbf{Model} &
\textbf{Size (KB)} & \textbf{INT8 / FP32 Acc.} & \textbf{Notes} \\
\midrule
\multirow{5}{*}{Vision}
& MobileNetV3 ($\times$1.0) & 2{,}478 & 93.56\% & 10-class, INT8 \\
& MobileNetV3 ($\times$0.75)& 1{,}495 & 93.11\% & 10-class, INT8 \\
& MobileNetV3 ($\times$0.5) & 703     & 91.67\% & 10-class, INT8 \\
& MobileNetV3 $80\times300$$^\ddagger$ & \textasciitilde1{,}200 &
  $(82.22 \pm 2.24)$\% &
  12-class, FP32 (s42: 83.17, s100: 79.67, s786: 83.83) \\
\midrule
\multirow{5}{*}{DS-CNN}
& 1a: Baseline     & 189 & 91.67\% & 10-class \\
& 1c: +SE          & 311 & 93.56\% & 10-class \\
& 1d: +Residual+MHSA & 302 & 94.22\% & 10-class \\
& 1e: All residuals  & 451 & 95.67\% & 10-class \\
& \textbf{v1: MynaNet} & \textbf{434} &
  \textbf{95.00\%$^\dagger$} & \textbf{10-class} \\
\midrule
TCN & 4c: Best TCN    & 450 & 92.44\% & 10-class \\
\bottomrule
\end{tabular}}
\end{table}

MynaNet~v1 achieves \SI{38}{\percent} smaller size than MobileNetV3-Small
($\times$0.5, 703~KB) while improving accuracy by +3.33\% on the 10-class
task. On the expanded 12-class dataset with identical $80\times300$
input and 80:10:10 split, MobileNetV3-Small achieves only
$82.22 \pm 2.24$\% FP32 (3-seed macOS), while the MCU-compatible Model~1c
achieves $93.52 \pm 0.58$\% FP32---a gap of more than \SI{11}{\percent}
absolute. This confirms that audio-specific inductive biases dominate
ImageNet pretraining for spectrogram classification at this dataset scale.

\subsection{Reproducibility: 3-Seed Evaluation}
\label{sec:reproducibility}

Table~\ref{tab:multiseed} reports the original 10-class Linux/CUDA results
for reference. Table~\ref{tab:multiseed-new} presents the re-validation on
the expanded 12-class dataset (macOS, Apple M4 Pro), using identical
hyperparameters (mels=64, mixup $\alpha=0.2$, 80:10:10 split).

\begin{table}[ht]
\centering
\caption{3-seed reproducibility evaluation (Linux, CUDA, 80:10:10 split,
         INT8 accuracy, 10-class dataset).}
\label{tab:multiseed}
\resizebox{\linewidth}{!}{%
\begin{tabular}{@{}lrrrll@{}}
\toprule
\textbf{Model} & \textbf{Seed 42} & \textbf{Seed 100} & \textbf{Seed 786} &
\textbf{Mean $\pm$ Std} & \textbf{Size} \\
\midrule
\textbf{MynaNet v1} & 94.00\% & 95.50\% & 95.50\% &
  \textbf{95.00\% $\pm$ 0.87\%} & \textbf{434~KB} \\
Model 1e & 94.33\% & 94.50\% & 95.17\% &
  94.67\% $\pm$ 0.45\% & 529~KB \\
\bottomrule
\end{tabular}
}
\end{table}

\begin{table}[ht]
\centering
\caption{3-seed re-validation on 12-class MyGardenBird dataset (macOS,
         Apple M4 Pro, 80:10:10 split, mels=64, mixup $\alpha=0.2$).
         FP32 and INT8 test accuracy reported.
         \TODO{Replace with Linux/CUDA re-run for authoritative figures.}}
\label{tab:multiseed-new}
\resizebox{\linewidth}{!}{%
\begin{tabular}{@{}lrrrrrl@{}}
\toprule
\textbf{Model} & \textbf{Seed 42} & \textbf{Seed 100} & \textbf{Seed 786} &
\textbf{Mean $\pm$ Std} & \textbf{Size} & \textbf{MCU} \\
\midrule
\multicolumn{7}{l}{\textit{FP32 accuracy}} \\
1c (DS-CNN+SE+Res) & 93.33\% & 94.17\% & 93.06\% &
  93.52\% $\pm$ 0.58\% & 311~KB & \checkmark \\
1d (+MHSA)         & 95.00\% & 95.14\% & 94.03\% &
  94.72\% $\pm$ 0.60\% & 302~KB & $\times$ \\
1e (+Wide)         & 93.47\% & 95.28\% & 94.44\% &
  94.40\% $\pm$ 0.91\% & 451~KB & $\times$ \\
\midrule
\multicolumn{7}{l}{\textit{INT8 accuracy}} \\
\textbf{1c (MCU-compatible)} & \textbf{93.19\%} & \textbf{94.31\%} & \textbf{93.06\%} &
  \textbf{93.52\% $\pm$ 0.69\%} & \textbf{311~KB} & \checkmark \\
1d (+MHSA)         & 95.00\% & 94.86\% & 94.03\% &
  94.63\% $\pm$ 0.52\% & 302~KB & $\times$ \\
1e (+Wide)         & 93.33\% & 95.14\% & 94.17\% &
  94.21\% $\pm$ 0.91\% & 451~KB & $\times$ \\
\bottomrule
\multicolumn{7}{l}{\footnotesize MCU: \checkmark = TFLite Micro compatible (Portenta H7);
  $\times$ = uses \texttt{BATCH\_MATMUL} (unsupported).} \\
\end{tabular}}
\end{table}

On the 12-class dataset, Model~1c (the only MCU-compatible architecture) achieves
93.52\%~$\pm$~0.69\% INT8 accuracy, a modest reduction of 0.48\% relative to its
10-class mean (94.00\%), consistent with the increased classification difficulty.
Model~1d leads on the expanded task at 94.63\%~$\pm$~0.52\% INT8, with 1e within
noise at 94.21\%~$\pm$~0.91\%. The higher variance of 1e (0.91\% vs 0.52\% for 1d)
suggests the wider channel configuration is less stable at this dataset scale.

% ==============================================================================
\section{Deployment Profile}
\label{sec:deployment}
% ==============================================================================

Table~\ref{tab:deployment} summarises the edge deployment characteristics
of MynaNet~v1 for an ARM Cortex-M7 processor at \SI{480}{MHz}.

\begin{table}[ht]
\centering
\caption{MynaNet v1 deployment profile.}
\label{tab:deployment}
\begin{tabular}{@{}lll@{}}
\toprule
\textbf{Metric} & \textbf{Value} & \textbf{Target / Notes} \\
\midrule
\multicolumn{3}{l}{\textit{Model size}} \\
\quad FP32 (.keras) & 1.58~MB & -- \\
\quad INT8 (.tflite) & 434~KB & Target: $<$512~KB \checkmark \\
\quad Margin under budget & 78~KB (15\%) & -- \\
\midrule
\multicolumn{3}{l}{\textit{Accuracy (10-class, Linux 3-seed)}} \\
\quad INT8 mean $\pm$ std & 95.00\% $\pm$ 0.87\% & Target: $>$90\% \checkmark \\
\quad INT8 min (seed 42)  & 94.00\% & $>$90\% \checkmark \\
\midrule
\multicolumn{3}{l}{\textit{Compute (estimated, Cortex-M7 @ 480 MHz)}} \\
\quad MACs             & 872{,}268   & -- \\
\quad Latency          & 116~ms      & Target: $<$500~ms \checkmark \\
\quad Throughput       & 8.6 inf/s   & Target: $>$2 inf/s \checkmark \\
\quad Energy / inference & 13.9~mJ  & -- \\
\midrule
\multicolumn{3}{l}{\textit{Training (Apple M4 Pro, 16~GB)}} \\
\quad Warmup (70 epochs)  & 126~min & -- \\
\quad Fine-tune (20 epochs) & 36~min & -- \\
\quad Total & 162~min & -- \\
\bottomrule
\end{tabular}
\end{table}

At \SI{116}{ms} latency and a \SI{10}{ms} frame stride, MynaNet~v1
processes a 3-second clip 24$\times$ faster than real-time, enabling
continuous monitoring with overlapping inference windows. Continuous
operation at 8.6 inferences per second draws approximately \SI{118}{mW},
feasible for solar-powered acoustic sensors.

\TODO{Add measured Cortex-M7 latency once firmware deployment is completed
on the ARGUS Smart Biologger prototype.  Current figure (116 ms) is
estimated from MAC count; actual measured latency may differ due to cache
effects and memory bandwidth.}


\begin{table}[t]
	\caption{Comparison of MynaNet against state-of-the-art architectures. Memory footprints represent quantized INT8 values where applicable.}
	\label{tab:comparison}
	\centering
	\footnotesize
	\setlength{\tabcolsep}{4pt} % Tightens horizontal spacing
	\renewcommand{\arraystretch}{1.1} % Slightly more vertical breathing room
	\begin{tabular}{l c c c c}
		\toprule
		\textbf{Model} & \textbf{Parameters} & \textbf{Flash} & \textbf{Peak SRAM} & \textbf{Accuracy} \\
		\midrule
		BirdNET \cite{kahl2021birdnet} & $\sim$27.0 M & $>$100 MB & N/A\textsuperscript{a} & \textbf{96.2\%} \\
		PANNs \cite{kong2020panns} & $\sim$81.0 M & 312 MB & N/A\textsuperscript{a} & 94.8\% \\
		MobileNetV3-S \cite{howard2019searching} & $\sim$2.5 M & 2.5 MB & $\sim$450 KB & 92.1\% \\
		DS-CNN (L) \cite{zhang2017hello} & 0.5 M & 500 KB & $\sim$180 KB & 89.4\% \\
		\midrule
		\textbf{MynaNet (Ours)} & \textbf{0.41 M} & \textbf{434 KB} & \textbf{112 KB} & \textbf{94.5\%} \\
		\bottomrule
		\multicolumn{5}{l}{\textsuperscript{a}\tiny Target device exceeds microcontroller-class memory (Inference on GPU/Cloud).}
	\end{tabular}
\end{table}

% ==============================================================================
\section{Discussion}
\label{sec:discussion}
% ==============================================================================

Our systematic exploration from MobileNetV3 baselines through TCN
experimentation to MynaNet~v1 yields nine transferable design insights.

\begin{enumerate}
  \item \textbf{Domain-specific inductive bias dominates raw capacity.}
        DS-CNN's factorised convolutions align with spectro-temporal
        structure, outperforming larger MobileNetV3 variants and 1D
        temporal models by exploiting 2D spatial locality.

  \item \textbf{Moderate attention is optimal.}
        Lightweight MHSA (2 heads, 88 dims) captures long-range temporal
        dependencies efficiently. Increasing to 3 heads (v2) provides no
        accuracy benefit at +43~KB cost, indicating attention-capacity
        saturation.

  \item \textbf{Residual connections scale depth.}
        Extending skip connections from 1 to all 4 blocks enabled 90-epoch
        training without degradation, unlocking the largest single gain
        (+1.45\%, 1d $\to$ 1e).

  \item \textbf{Less can be more.}
        The 6\% channel reduction in v1 ([120, 180, 146, 360]) eliminated
        over-parameterisation, improving accuracy +2.00\% over v0 while
        reducing size by 47~KB. The production-optimised v1 surpasses the
        larger 1e in authoritative multi-seed evaluation.

  \item \textbf{Quantisation-friendly design.}
        ReLU6 activations and batch normalisation minimise PTQ accuracy
        drop to $<$0.5\%, meeting deployment constraints without
        quantisation-aware training overhead.

  \item \textbf{Mixup $>$ SpecAugment for bird calls.}
        With 4{,}800 training samples (10-class dataset), Mixup's label smoothing generalises
        better (+0.83\%) than SpecAugment's masking, which risks erasing
        discriminative short vocalisations.

  \item \textbf{Spectrogram resolution is critical.}
        80~mel bins (100~Hz/bin) vs.\ 64~bins (125~Hz/bin) improves
        high-frequency species discrimination, particularly for species
        with overlapping pitch ranges.

  \item \textbf{TCNs are ill-suited for spectrogram classification.}
        The 511-frame receptive field (5.1~s) exceeds clip duration;
        causal padding incurs unnecessary computation; and 1D operation
        discards 2D spectral structure. TCNs may be preferable for
        streaming keyword spotting over full-clip classification.

  \item \textbf{Training data volume matters at this scale.}
        An 80:10:10 split (4{,}800 train) outperforms 70:15:15 (4{,}200
        train) by +1.56\% on the 10-class dataset, showing that additional
        training samples outweigh a larger test set at this dataset scale.
\end{enumerate}

\TODO{Expand Discussion with: (1) comparison to BirdNET and other published
bioacoustics models on a common species/dataset; (2) limitations ---
dataset covers only 12 Southeast Asian species, performance on unseen
species / environments unknown; (3) future work: knowledge distillation from
EfficientNet-B0/ResNet-50 ensembles, learnable mel
filterbanks~\cite{zeghidour2018learning}, NAS for channel-width
optimisation, multi-task joint classification + call-type detection.}

% ==============================================================================
\section{Conclusion}
\label{sec:conclusion}
% ==============================================================================

We presented MynaNet~v1, a 434~KB INT8 CNN achieving 95.00\% $\pm$ 0.87\%
mean accuracy on a ten-class bird call dataset, and re-validated on an
expanded twelve-class dataset where the MCU-compatible Model~1c achieves
93.52\% $\pm$ 0.69\% INT8 accuracy. Through 20+ controlled experiments
spanning MobileNetV3 baselines, TCN sequence models, and audio-specific
DS-CNN variants, we demonstrated that domain-aware architecture design
substantially outperforms generic solutions within extreme edge deployment
constraints. The key contributions are:

\begin{itemize}
  \item A principled architecture evolution from MobileNetV3 (91.67\%,
        703~KB) to MynaNet~v1 (95.00\%, 434~KB), achieving +3.33\% accuracy
        with 38\% smaller footprint.
  \item Comprehensive ablation studies isolating the contribution of
        depthwise separation, SE attention, residual-all connectivity,
        lightweight MHSA, Mixup augmentation, and INT8 quantisation.
  \item A deployment-ready model meeting all edge constraints:
        $<$512~KB, $>$90\% accuracy, $<$500~ms latency on Cortex-M7.
  \item Empirical evidence that 2D audio-specific convolutions outperform
        1D temporal models by 3.2\% for spectrogram-based classification.
\end{itemize}

MynaNet~v1 demonstrates that careful, principled capacity allocation ---
rather than naive parameter scaling --- is the key to production-grade
accuracy within the extreme constraints of wildlife acoustic monitoring
hardware.

% ==============================================================================
\section*{Data Availability}
\label{sec:data_availability}
% ==============================================================================

The MyGardenBird dataset associated with this study is available in the Zenodo repository, \cite{zabidi2026dataset}. The dataset includes annotated audio recordings and the specific train/validation/test partitions used for the MynaNet benchmarking experiments.

\section*{Declaration of Competing Interest}
The authors declare that they have no known competing financial interests or personal relationships that could have appeared to influence the work reported in this paper.

\section*{Acknowledgments}
The authors would like to thank [Grant Body/University] for supporting this research under grant [Grant Number]. We also express our gratitude to the citizen scientists of the Xeno-canto community whose recordings contributed to the foundational understanding of the species studied.



% ==============================================================================
% References
% ==============================================================================

\bibliographystyle{elsarticle-num}
\bibliography{mynanet_neurocomputing}

\end{document}
